\documentclass[12pt]{article}
\usepackage[margin=1in]{geometry}
\usepackage{amsmath}

\begin{document}

\noindent
\textbf{Name:} \underline{\hspace{6cm}} \hfill \textbf{Date:} \underline{\hspace{4cm}} \\
\textbf{Course:} CS 471 - Artificial Intelligence \\
\textbf{Lesson 15:} Reinforcement Learning III (Q-Learning)

\vspace{1cm}

\section*{Section 1: Instructor-Led Walkthrough}
\textbf{Objective:} Calculate a Q-Value update for an autonomous perimeter sentry. \\

Given the Q-learning update equation:
\[ Q_{k+1}(s,a) \leftarrow (1-\alpha)Q_k(s,a) + \alpha \left[ R(s,a,s') + \gamma \max_{a'} Q_k(s',a') \right] \]

\noindent
Assume the sentry is at Checkpoint 1 ($s$), takes action "Advance" ($a$), and receives a reward of $R = -2$. It transitions to Checkpoint 2 ($s'$), which has a maximum future Q-value of $10$. \\
Let $\alpha = 0.5$, $\gamma = 1.0$, and the initial $Q(s,a) = 0$. \\

\noindent
\textbf{Step 1: Calculate the sample} \\
$sample = $

\vspace{2cm}

\noindent
\textbf{Step 2: Calculate the new Q-value} \\
$Q(s,a) \leftarrow $

\vspace{2cm}

\section*{Section 2: Student Practice}

\textbf{Problem 1: Exploration vs. Exploitation} \\
An Electronic Warfare algorithm utilizes an $\epsilon$-greedy strategy to find optimal frequencies for jamming adversary comms. If $\epsilon = 0.1$, explain the behavioral tradeoff the algorithm makes at each time step.

\vspace{4cm}

\textbf{Problem 2: Model-Free Learning} \\
Why does an autonomous drone using Q-Learning not require prior knowledge of the transition function $T(s,a,s')$ to learn an optimal flight path? 

\vspace{4cm}

\end{document}